\documentclass{amestyle}

%%%%%%%%%%%%%%%%%%%%%%%%%%%%%%%%%%%%%%%%%%%%%%%%%%%%%%%%%%%%%%%%%%%%%%%%%%%%%%
%                         MANDATORY PACKAGES                                 %
%%%%%%%%%%%%%%%%%%%%%%%%%%%%%%%%%%%%%%%%%%%%%%%%%%%%%%%%%%%%%%%%%%%%%%%%%%%%%%
\usepackage[polish,english]{babel}
\usepackage[T1]{fontenc}
\usepackage[utf8]{inputenc}        % Please use UTF-8 encoding
\usepackage[numbers,sort&compress]{natbib}
\usepackage{courier}
\usepackage{tgtermes,newtxtext,newtxmath} % Times fonts
\usepackage{marvosym} % for graphics (Letter)

%%%%%%%%%%%%%%%%%%%%%%%%%%%%%%%%%%%%%%%%%%%%%%%%%%%%%%%%%%%%%%%%%%%%%%%%%%%%%%
%                         MANDATORY SETTINGS                                 %
%%%%%%%%%%%%%%%%%%%%%%%%%%%%%%%%%%%%%%%%%%%%%%%%%%%%%%%%%%%%%%%%%%%%%%%%%%%%%%
\setlength{\bibsep}{0pt plus 4pt}
\renewcommand{\bibfont}{\footnotesize}
\everymath{\displaystyle}


%%%%%%%%%%%%%%%%%%%%%%%%%%%%%%%%%%%%%%%%%%%%%%%%%%%%%%%%%%%%%%%%%%%%%%%%%%%%%%
%                         OPTIONAL PACKAGES                                  %
%%%%%%%%%%%%%%%%%%%%%%%%%%%%%%%%%%%%%%%%%%%%%%%%%%%%%%%%%%%%%%%%%%%%%%%%%%%%%%
\usepackage{amsmath,amsfonts}      % Note: amsmath should go before hyperref
\usepackage[pdftex]{graphicx}
\usepackage[pdftex,colorlinks,allcolors=blue]{hyperref}
\usepackage{listings}
\usepackage{tikz}
\usepackage{subcaption}
\usepackage{lipsum}


%%%%%%%%%%%%%%%%%%%%%%%%%%%%%%%%%%%%%%%%%%%%%%%%%%%%%%%%%%%%%%%%%%%%%%%%%%%%%%
%                                                                            %
%                                                                            %
%               METADATA TO BE PROVIDED BY AME EDITOR                        %
%           (Authors should leave this section unaltered)                    %
%                                                                            %
%                                                                            %
%%%%%%%%%%%%%%%%%%%%%%%%%%%%%%%%%%%%%%%%%%%%%%%%%%%%%%%%%%%%%%%%%%%%%%%%%%%%%%
%
%% \amevolume{66}                  % Volume number
%% \ameissue{1}                       % Issue number
%% \ameyear{2019}                     % Year
%% \ameprintpage{49}                  % First page of the paper in the book
%% \amedoi{10.24425/ame.2019.12345}  % DOI number
%% \amereceived{September 09, 2018}   % Date of the manuscript reception
%% \ameaccepted{February 22, 2019}    % Date of the final version
%
%%%%%%%%%%%%%%%%%%%%%%%%%%%%%%%%%%%%%%%%%%%%%%%%%%%%%%%%%%%%%%%%%%%%%%%%%%%%%


%%%%%%%%%%%%%%%%%%%%%%%%%%%%%%%%%%%%%%%%%%%%%%%%%%%%%%%%%%%%%%%%%%%%%%%%%%%%%%
%                                                                            %
%                                                                            %
%                 METADATA TO BE FILLED IN BY AUTHORS                        %
%              (Please provide necessary information below)                  %
%                                                                            %
%%%%%%%%%%%%%%%%%%%%%%%%%%%%%%%%%%%%%%%%%%%%%%%%%%%%%%%%%%%%%%%%%%%%%%%%%%%%%%

%%%%%%%%%%%%%%%%%%%%%%%%%%%%%%%%%%%%%%%%%%%%%%%%%%%%%%%%%%%%%%%%%%%%%%%%%%%%%%
% AUTHORS: Use \title command to provide full title of the paper.
%%%%%%%%%%%%%%%%%%%%%%%%%%%%%%%%%%%%%%%%%%%%%%%%%%%%%%%%%%%%%%%%%%%%%%%%%%%%%%
\title{\LaTeX\ template for authors preparing manuscripts to the Archive of
       Mechanical Engineering}

%%%%%%%%%%%%%%%%%%%%%%%%%%%%%%%%%%%%%%%%%%%%%%%%%%%%%%%%%%%%%%%%%%%%%%%%%%%%%%
% AUTHORS: Use \shorttitle to customize short title. This title will appear in
%          running headers and possibly in some other places, where short title
%          is desirable. Note, that the title used specifically for running
%          head may also be defined with \headtitle (below).
%%%%%%%%%%%%%%%%%%%%%%%%%%%%%%%%%%%%%%%%%%%%%%%%%%%%%%%%%%%%%%%%%%%%%%%%%%%%%%
%\shorttitle{Short title of the manuscript}

%%%%%%%%%%%%%%%%%%%%%%%%%%%%%%%%%%%%%%%%%%%%%%%%%%%%%%%%%%%%%%%%%%%%%%%%%%%%%%
% AUTHORS: \headtitle may be used to customize the title in running headers.
%%%%%%%%%%%%%%%%%%%%%%%%%%%%%%%%%%%%%%%%%%%%%%%%%%%%%%%%%%%%%%%%%%%%%%%%%%%%%%
% \headtitle{Title in running header}

%%%%%%%%%%%%%%%%%%%%%%%%%%%%%%%%%%%%%%%%%%%%%%%%%%%%%%%%%%%%%%%%%%%%%%%%%%%%%%
% AUTHORS: \author command is provided to define authors' names and
%          affiliations and identify corresponding author. At the beginning 
%          corresponding author name and email should be given using 
%          command \corrauthor. 
%          Author records should be separated with the \and command. 
%          Each record should contain at least author's name. Author's
%          name may be followed with \contact{...} command providing contact
%          information. If two or more authors have same affiliation, the
%          affiliation should be defined after the first author's name and then
%          \contactmark[N] may be used for others authors to refer to the
%          affiliation record of the previous author. The number N is the
%          number of the author's affifliation being referred to (see example below).
%          \orcidid command is optional and the ORCID icon redirects reader to author's
%          account on orcid.org
%%%%%%%%%%%%%%%%%%%%%%%%%%%%%%%%%%%%%%%%%%%%%%%%%%%%%%%%%%%%%%%%%%%%%%%%%%%%%%
\author{%
\corrauthor{Second Author, email:  \email{corresponding@inst.edu}}%
%special contact to indicate corresponding author - put his/her name here [1]
  First Author\orcidid{0000-0002-1825-0097}% ORCID number is optional
    \contact{First Author's Institute, city, country. % [2]
             Emails: \email{first.author@institute1.edu},\newline % e-mail is optional
                        \email{third.author@institute1.edu} % e-mail is optional
                }%
  \ \ \and
  Second Author\orcidid{0000-0002-1825-0097}% ORCID number is optional 
    \contact{Second and Third Author's Institute, city, country.}%
  \ \ \and
  Third Author% 
    \contactmark[2]% this refers to the affiliation [2] above
}

%%%%%%%%%%%%%%%%%%%%%%%%%%%%%%%%%%%%%%%%%%%%%%%%%%%%%%%%%%%%%%%%%%%%%%%%%%%%%%
% AUTHORS: \headauthor should be used to customize author information appearing in
%          the running headers (to omit the orcid icon connected with authors names).
%%%%%%%%%%%%%%%%%%%%%%%%%%%%%%%%%%%%%%%%%%%%%%%%%%%%%%%%%%%%%%%%%%%%%%%%%%%%%%
\headauthor{F. Author \and S. Author}

%%%%%%%%%%%%%%%%%%%%%%%%%%%%%%%%%%%%%%%%%%%%%%%%%%%%%%%%%%%%%%%%%%%%%%%%%%%%%%
% AUTHORS: \keywords should be used to provide a list of key words separated
%          with comma.
%%%%%%%%%%%%%%%%%%%%%%%%%%%%%%%%%%%%%%%%%%%%%%%%%%%%%%%%%%%%%%%%%%%%%%%%%%%%%%
\keywords{First keyword \and second keyword \and third keyword}

%%%%%%%%%%%%%%%%%%%%%%%%%%%%%%%%%%%%%%%%%%%%%%%%%%%%%%%%%%%%%%%%%%%%%%%%%%%%%%
%                           LATEX SETTINGS                                   %
%%%%%%%%%%%%%%%%%%%%%%%%%%%%%%%%%%%%%%%%%%%%%%%%%%%%%%%%%%%%%%%%%%%%%%%%%%%%%%
\selectlanguage{english}


\begin{document}

\maketitle

\license

%%%%%%%%%%%%%%%%%%%%%%%%%%%%%%%%%%%%%%%%%%%%%%%%%%%%%%%%%%%%%%%%%%%%%%%%%%%%%%
\begin{abstract}
  \label{abstract}
  \lipsum[1]
\end{abstract}
%%%%%%%%%%%%%%%%%%%%%%%%%%%%%%%%%%%%%%%%%%%%%%%%%%%%%%%%%%%%%%%%%%%%%%%%%%%%%%





%%%%%%%%%%%%%%%%%%%%%%%%%%%%%%%%%%%%%%%%%%%%%%%%%%%%%%%%%%%%%%%%%%%%%%%%%%%%%%
\section{Introduction}
\label{sec:intro}

\lipsum[1]

And this paragraph shows how we cite multiple positions at once. The citations
shall get compressed automatically
\cite{smith.kowalski:running,%
      anderson.duan:2000:highly,%
      arnold:2007:multi,%
      arveson:2001:spectral:book}.
Second trial should lead to compressed citations too
\cite{smith.kowalski:running,%
      anderson.duan:2000:highly,%
      andrus:1979:numerical,%
      arnold:1992:ordinary,%
      arnold:2007:multi,%
      arnold.gunther:2001:preconditioned}.

The organization of this paper is following. In section \ref{sec:first} we have
Lorem Ipsum itemized and enumerated lists. In section \ref{sec:first-first}
we have equations. In section \ref{sec:first-second} we have table and figures.
Section \ref{sec:conclusions} concludes the paper.


%%%%%%%%%%%%%%%%%%%%%%%%%%%%%%%%%%%%%%%%%%%%%%%%%%%%%%%%%%%%%%%%%%%%%%%%%%%%%%
\section{First section}
\label{sec:first}

\lipsum[2]
%%%%%%%%%%%%%%%%%%%%%%%%%%%%%%
\begin{itemize}
  \item Lorem ipsum dolor sit amet,
  \item vestibulum ut, placerat ac,
  \item adipiscing vitae, felis.
\end{itemize}
%%%%%%%%%%%%%%%%%%%%%%%%%%%%%%

\lipsum[3]
%%%%%%%%%%%%%%%%%%%%%%%%%%%%%%
\begin{enumerate}
  \item Lorem ipsum dolor sit amet,
  \item vestibulum ut, placerat ac,
  \item adipiscing vitae, felis.
\end{enumerate}
%%%%%%%%%%%%%%%%%%%%%%%%%%%%%%

\lipsum[4]

%%%%%%%%%%%%%%%%%%%%%%%%%%%%%%%%%%%%%%%%%%%%%%%%%%%%%%%%%%%%%%%%%%%%%%%%%%%%%%
\subsection{First subsection of first section}
\label{sec:first-first}

\lipsum[5]
%%%%%%%%%%%%%%%%%%%%%%%%%%%%%%
\begin{equation}
  \mathbf{A} = \begin{bmatrix}
    \boldsymbol{\lambdaup} & \boldsymbol{\Phi} \\
    \mathbf{B} & \boldsymbol{\muup}
  \end{bmatrix}
  \label{eq:1}
\end{equation}
%%%%%%%%%%%%%%%%%%%%%%%%%%%%%%
We may always refer the equation \eqref{eq:1} with \verb|\eqref|.

\lipsum[6]
%%%%%%%%%%%%%%%%%%%%%%%%%%%%%%
\begin{subequations}
\label{eq:2}
\begin{align}
  a & = b + c \label{eq:2a}
  \\
  d & = e + f \label{eq:2b}
\end{align}
\end{subequations}
%%%%%%%%%%%%%%%%%%%%%%%%%%%%%%

\lipsum[7]

%%%%%%%%%%%%%%%%%%%%%%%%%%%%%%%%%%%%%%%%%%%%%%%%%%%%%%%%%%%%%%%%%%%%%%%%%%%%%%
\subsection{Second subsection of second section}
\label{sec:first-second}

\lipsum[8]
%%%%%%%%%%%%%%%%%%%%%%%%%%%%%%
\begin{table}[htbp]
  \caption{Number of cars produced in the world}
  \label{tab:1}
  \centering
  \begin{tabular}{|l|r|}
    \hline
    \multicolumn{1}{|c|}{Year} &
    \multicolumn{1}{|c|}{Cars produced} \\\hline
    2009 & 47~772~598 \\
    2010 & 58~264~852 \\
    2011 & 59~929~016 \\\hline
  \end{tabular}
\end{table}
%%%%%%%%%%%%%%%%%%%%%%%%%%%%%%
We can always refer the Table~\ref{tab:1} with \verb|\ref|. Below shall also
be a~figure. We can always refer the Fig.~\ref{fig:fbar} with \verb|\ref|.
%%%%%%%%%%%%%%%%%%%%%%%%%%%%%%
\begin{figure}[htbp]
  \centering
  \includegraphics[width=0.5\textwidth]{fbar.pdf}
  \caption{Four-bar mechanism with all coordinates defined}
  \label{fig:fbar}
\end{figure}
%%%%%%%%%%%%%%%%%%%%%%%%%%%%%%

\lipsum[10]
%%%%%%%%%%%%%%%%%%%%%%%%%%%%%%
\begin{figure}[htbp]
  \centering
  \begin{tikzpicture}[%
      sibling distance  = 10em,%
      every node/.style = {%
        shape=rectangle, rounded corners, draw,%
        align=center, top color=white,%
        bottom color=blue!20}]
    \node {Mechanics}
      child { node {Kinematics} }
      child { node {Dynamics} };
  \end{tikzpicture}
  \caption{A figure created with TikZ commands}
  \label{fig:tikz}
\end{figure}
%%%%%%%%%%%%%%%%%%%%%%%%%%%%%%

\lipsum[11]
%%%%%%%%%%%%%%%%%%%%%%%%%%%%%%
\begin{figure}[htbp]
  \centering
  \begin{subfigure}[b]{0.3\textwidth}
    \includegraphics[width=\textwidth]{fbar.pdf}
    \caption{Initial position}
    \label{fig:fbar2a}
  \end{subfigure}
  \begin{subfigure}[b]{0.3\textwidth}
    \includegraphics[width=\textwidth]{fbar.pdf}
    \caption{Final position}
    \label{fig:fbar2b}
  \end{subfigure}
  \caption{Four bar in its Initial and Final position}
  \label{fig:fbar2}
\end{figure}
%%%%%%%%%%%%%%%%%%%%%%%%%%%%%%



%%%%%%%%%%%%%%%%%%%%%%%%%%%%%%%%%%%%%%%%%%%%%%%%%%%%%%%%%%%%%%%%%%%%%%%%%%%%%%
\section{Conclusions}
\label{sec:conclusions}

\lipsum[12]

%%%%%%%%%%%%%%%%%%%%%%%%%%%%%%%%%%%%%%%%%%%%%%%%%%%%%%%%%%%%%%%%%%%%%%%%%%%%%%
%                             APPENDICES                                     %
%%%%%%%%%%%%%%%%%%%%%%%%%%%%%%%%%%%%%%%%%%%%%%%%%%%%%%%%%%%%%%%%%%%%%%%%%%%%%%
\begin{appendix}
  %%%%%%%%%%%%%%%%%%%%%%%%%%%%%%%%%%%%%%%%%%%%%%%%%%%%%%%%%%%%%%%%%%%%%%%%%%%%
  \section{Appendix A}
  \label{sec:appendix-a}
  \lipsum[1]

  %%%%%%%%%%%%%%%%%%%%%%%%%%%%%%%%%%%%%%%%%%%%%%%%%%%%%%%%%%%%%%%%%%%%%%%%%%%%
  \section{Appendix B}
  \label{sec:appendix-b}
  \lipsum[2]
\end{appendix}

\begin{acknowledgements}
  Thank you, thank you, thank you. Thank you, thank you, thank you. Thank you,
  thank you, thank you.  Thank you, thank you, thank you. Thank you, thank you,
  thank you. Thank you, thank you, thank you. Thank you, thank you, thank you.
  Thank you, thank you, thank you.
\end{acknowledgements}

%%%%%%%%%%%%%%%%%%%%%%%%%%%%%%%%%%%%%%%%%%%%%%%%%%%%%%%%%%%%%%%%%%%%%%%%%%%%%%
%                      LEAVE THE FOLLOWING LINES UNALTERED                   %
%%%%%%%%%%%%%%%%%%%%%%%%%%%%%%%%%%%%%%%%%%%%%%%%%%%%%%%%%%%%%%%%%%%%%%%%%%%%%%
\begin{receiptnote}
  Manuscript received by Editorial Board, \theamereceived; \\
  final version, \theameaccepted.
\end{receiptnote}



%%%%%%%%%%%%%%%%%%%%%%%%%%%%%%%%%%%%%%%%%%%%%%%%%%%%%%%%%%%%%%%%%%%%%%%%%%%%%%
%                             REFERENCES                                     %
%%%%%%%%%%%%%%%%%%%%%%%%%%%%%%%%%%%%%%%%%%%%%%%%%%%%%%%%%%%%%%%%%%%%%%%%%%%%%%
\bibliographystyle{unsrt}
\bibliography{References}


\end{document}
% vim: set spell expandtab tabstop=2 shiftwidth=2 syntax=tex ff=unix:
